\section{The Impediment and the Penal Backstop}

Canon 277 §1 states an obligation, and an obligation on its own is a moral fact with no juridical grip. The Code gives it two kinds of grip: it makes the cleric incapable of contracting marriage, and it makes the attempt a delict. These are different instruments and they do different work, which is why the Code puts them four hundred canons apart.

\subsection{Canon 1087: incapacity, not prohibition}

\latin{Invalide matrimonium attentant, qui in sacris ordinibus sunt constituti}---those who are constituted in sacred orders invalidly attempt marriage. Six words of Latin, and the whole of the Latin Church's law on the point.\endnote{CIC c.~1087, Latin and English deliveries, fetched, hashed, and read 2026-07-25. The formulation is exactly parallel to the Eastern Code's c.~804 (section~7), and both use \latin{attentare}---to attempt---which is the technical verb for an act that has the outward form of the juridical act and none of its effect.}

The canon does not forbid; it disables. A diriment impediment (\latin{impedimentum dirimens}) renders a person incapable of validly contracting, so that the marriage a cleric attempts is not an illicit marriage but not a marriage. The distinction has consequences that a prohibition would not produce: there is nothing to dissolve, no annulment process is required to establish that nothing happened, the civil act has no canonical effect whatever, and any children of the union are affected in canon law only through the norms on legitimacy and not through anything in this canon.

Three features of the canon's reach are worth stating exactly. It attaches to \emph{sacred orders}, that is, to the diaconate, the presbyterate, and the episcopate alike---so a permanent deacon who is widowed falls under it as fully as a bishop does. It attaches to the \emph{order}, not to the obligation of celibacy, which is why it binds a married permanent deacon after his wife's death even though he never assumed the obligation of canon 1037. And it is an impediment of \emph{ecclesiastical} law, not of divine or natural law, which is what makes the next paragraph possible at all.

\subsection{Who can dispense it, and the one striking exception}

Because the impediment is ecclesiastical, the Church can relax it, and the Code says exactly who may. Canon 1078 §1 gives the local ordinary a general power to dispense his subjects from impediments of ecclesiastical law, \latin{exceptis iis, quorum dispensatio Sedi Apostolicae reservatur}; §2 1° then reserves ``the impediment arising from sacred orders or from a public perpetual vow of chastity in a religious institute of pontifical right.'' The ordinary route, then, is a reserved dispensation from the Apostolic See, which in practice is sought and granted together with the rescript of loss of the clerical state (section~6).

The exception is canon 1079 §1, and it is a genuine surprise. In urgent danger of death the local ordinary may dispense his subjects from the form of celebration and from ``each and every impediment of ecclesiastical law, whether public or occult,'' \latin{excepto impedimento orto ex sacro ordine presbyteratus}---except the impediment arising from the sacred order of \emph{presbyterate}. The reservation is to the presbyterate only. A deacon in danger of death may therefore be dispensed by his own local ordinary---and, where the ordinary cannot be reached, by his pastor, a properly delegated sacred minister, or the priest or deacon assisting under canon 1116 §2 (§2), and by a confessor for occult impediments in the internal forum (§3).\endnote{CIC cc.~1078, 1079, 1080, Latin and English deliveries, fetched, hashed, and read 2026-07-25. Canon 1080 §1 supplies the second urgent case---the impediment discovered when everything is prepared for the wedding---and there the exception is wider: it excludes precisely the impediments of c.~1078 §2, so the sacred-orders impediment at any grade is reserved. The two urgency canons therefore differ, and the difference is not a slip: the danger-of-death canon narrows the reservation to the presbyterate; the everything-is-ready canon does not. The Eastern Code draws the same narrow line at CCEO c.~796 §1, excepting \latin{impedimentum ordinis sacri sacerdotii}.} That is a small, precise, and easily missed provision, and it demonstrates something about the shape of the whole discipline: the Church treats the diaconal and the presbyteral impediment differently even inside a canon most readers will never have occasion to use.

\subsection{Canon 1394 §1: the delict}

The penal law approaches the same facts from the other side. In Book~VI as integrally revised in 2021: \latin{Clericus matrimonium, etiam civiliter tantum, attentans, in suspensionem latae sententiae incurrit, firmis praescriptis cann.~194, §~1, n.~3, et 694, §~1, n.~2; quod si monitus non resipuerit vel scandalum dare perrexerit, gradatim privationibus vel etiam dimissione e statu clericali puniri debet}---a cleric who attempts marriage, ``even if only civilly,'' incurs a \latin{latae sententiae} suspension; if after warning he has not reformed or continues to give scandal, ``he must be progressively punished by deprivations, or even by dismissal from the clerical state.''\endnote{CIC c.~1394 §1 in the text of Book~VI as revised by \emph{Pascite gregem Dei} (23 May 2021, in force 8 December 2021); Latin delivery of Book~VI and English delivery of cann.~1364--1399 fetched, hashed, and read 2026-07-25, both byte-identical to the states registered in this repository's source library. Section~2 of the canon reaches a perpetually professed religious who is not a cleric. The cross-references are to c.~194 §1 3° (removal from ecclesiastical office by the law itself for a cleric who has attempted marriage even civilly) and c.~694 §1 2° (a religious who has attempted marriage even civilly is to be held \latin{ipso facto} dismissed from the institute).}

Four points about this canon are commonly muddled. The phrase \latin{etiam civiliter tantum} exists because the canonical marriage was impossible: the delict consists in the attempt, and a civil ceremony is an attempt. The suspension is \latin{latae sententiae}, incurred by the act itself without any sentence, and the 2021 revision retained it in a book that reduced automatic penalties elsewhere. The escalation clause is not automatic: it requires a warning, and unrepentance or continuing scandal, and it is \emph{graduated}. And dismissal from the clerical state under this canon is a penalty, imposed for a delict; it is not the same act as the rescript of dispensation that a cleric petitions for. Both end in the loss of the clerical state, and one is a punishment while the other is a favour---which is the subject of the next section.

\subsection{The neighbouring canons, and one boundary the legislator drew}

Canon 1394 does not stand alone. Canon 1395 §1 punishes with suspension the cleric living in concubinage ``other than in the case mentioned in can.~1394'' and the cleric who ``continues in some other external sin against the sixth commandment of the Decalogue which causes scandal,'' with graduated additional penalties up to dismissal if he persists after warning; §2 reaches other public offences against the sixth commandment; §3 reaches offences committed by force, threats, or abuse of authority, and the coercion of anyone to perform or submit to sexual acts.

Where the revised Book~VI puts each of these is itself a legislative judgment worth recording. Canons 1392--1396, including both the attempted marriage and the concubinage, stand under Title~V, \emph{Offences against special obligations}. The delicts of canon 1398---offences against the sixth commandment with a minor or with a person who habitually has an imperfect use of reason or to whom the law recognizes equal protection, and the associated grooming and pornography delicts---stand under Title~VI, \emph{Offences against human life, dignity and freedom}, alongside homicide, abduction, mutilation, and abortion.\endnote{CIC Book~VI, Part~II, Titles~V and~VI (cc.~1392--1398), Latin and English deliveries fetched, hashed, and read 2026-07-25. The 2021 revision's relocation of these delicts is a matter of the promulgated text's own structure, which is all that is claimed here; the drafting history was not examined. What follows from the structure is limited but real: the legislator classifies the abuse of a minor as an offence against that person, not as a failure of a cleric's special obligations.} The legislator has thereby declined to classify the abuse of a minor as a failure of clerical continence. It is classified as a crime against the person harmed. Section~13 returns to what that does and does not settle.
